\section{Регулярные выражения}

Регулярные выражения~--- один из часто применяющихся на практике классических способов задать регулярный язык%
\sidenote{
    Замечание для программистов.
    Важно понимать, что речь идёт о формальной конструкции, а не о том, что называется регулярными выражениями в различных языках программирования, библиотеках, инструментах, где под названием \enquote{регулярные выражения} могут скрываться конструкции, существенно более выразительные, чем обсуждаемые здесь.
}.
Основывается этот способ на предложении синтаксиса для описания \emph{регулярных множеств}%
\sidenote{Помним, что язык~--- это множество слов.}.

\begin{definition}[Регулярное множество]
    Регулярное множество (над алфавитом $\Sigma$) это:
    \begin{itemize}
        \item $\varnothing$ или пустое множество;
        \item $\{\varepsilon\}$ или множество, состоящее только из пустой цепочки;
        \item $\{t\}$, $t \in \Sigma$ или множество, состоящее из одного символа из алфавита;
        \item $R_1 \cup R_2$ или объединение множеств, где $R_1$ и $R_2$~--- регулярные множества;
        \item $R_1 \cdot R_2$ или конкатенация множеств, где $R_1$ и $R_2$~--- регулярные множества;
        \item $R^*$ или замыкание Клини, где $R$~--- регулярное множество.
    \end{itemize}
\end{definition}

Для того, чтобы описывать такие множества, можно пользоваться \emph{регулярными выражениями}, которые, фактически, предоставляют синтаксис для описания операций над регулярными множествами, представленными выше.

\begin{definition}[Регулярное выражение]
    Регулярное выражение (над алфавитом $\Sigma$) это:
    \begin{itemize}
        \item пустое множество $\varnothing$;
        \item множество из пустой цепочки $\varepsilon$;
        \item множество из одного символа $t$, $t \in \Sigma$;
        \item объединение регулярных выражений (множеств) $R_1 \mid R_2$, где $R_1$ и $R_2$~--- регулярные выражения;
        \item конкатенация регулярных выражений (множеств) $R_1 \cdot R_2$, где $R_1$ и $R_2$~--- регулярные выражения;
        \item замыкание клини $R^*$, где $R$~--- регулярное выражение;
        \item группирующие скобки $(R)$, где $R$~--- регулярное выражение.
    \end{itemize}
\end{definition}

Отметим несколько важных с прикладной точки зрения моментов.
Во-первых, часто используется расширенный синтаксис, в который добавляются конструкции не увеличивающие выразительную силу, но упрощающие запись.
Например, встречаются следующие расширения%
\sidenote{
    Существуют и другие, однако их мы не будем использовать и, соответственно, рассматривать.
    Читатель может вспомнить, что называется регулярными выражениями в его любимом языке программирования и попробовать самостоятельно выразить имеющиеся там конструкции через базовые.}:
\begin{itemize}
    \item $R? = R \mid \varepsilon$, где $R$~--- регулярное выражение;
    \item $R^+ = R \cdot R^*$, где $R$~--- регулярное выражение.
\end{itemize}

Во-вторых, конструкции $\varnothing$ и $\varepsilon$ используются крайне редко, особенно в случае расширенного синтаксиса, так как часто выражение, эквивалентное использующему данные конструкции, более компактно записывается с использованием расширенного синтаксиса.
В-третьих, оператор конкатенации часто опускается%
\sidenote{Как и знак умножения во многих математических записях.}.

Рассмотрим несколько примеров регулярных выражений.
\begin{example}
    Регулярное выражение $a$ задаёт регулярное множество $\{a\}$ и, соответственно, язык из единственного слова $a$.
\end{example}


\begin{example}
    Регулярное выражение $ab$ задаёт регулярное множество $\{ab\}$ и, соответственно, язык из единственного слова $ab$.
\end{example}


\begin{example}\label{exmpl:regexp_a_star}
    Регулярное выражение $a^*$ задаёт регулярное множество \[R = \bigcup_{i=0}^{\infty}{a^i} = \{\varepsilon, a, aa, aaa, \ldots \}\] и, соответственно, бесконечный язык, содержащий для любого неотрицательного целого $n$ цепочку из символов $a$ длины $n$.
\end{example}

\begin{example}
    Регулярное выражение $a^*b$ задаёт бесконечный язык, содержащий цепочки, полученные конкатенацией цепочек из языка, описанного в примере~\ref{exmpl:regexp_a_star} и символа $b$.
    Иными словами, любая цепочка из языка выглядит как последовательность символов $a$ произвольной длины (возможно пустая), за которой следует символ $b$.
    Примеры цепочек из языка: $b, ab, aaab, aaaaab$.
\end{example}

\begin{example}
    Регулярное выражение $(a\mid b)^*$ задаёт бесконечный язык, цепочки которого имеют произвольную длину и состоят из символов $a$ и $b$ в произвольном порядке.
    Примеры цепочек из языка: $\varepsilon, a, aa, b, ba, aba$.
\end{example}

\begin{example}
    Регулярное выражение\sidenote{Здесь мы воспользовались расширенным синтаксисом.} $(ab)^*c?$ задаёт бесконечный язык, цепочки которого состоят их произвольного количества повторений подпоследовательности $ab$ за которой может следовать (но не обязательно следует) символ $c$.
    Примеры цепочек из языка: $\varepsilon, abab, abc, abababc$.
\end{example}
